\documentclass[12pt,a4paper]{article}

% ── Packages ──────────────────────────────────────────────────────────────────
\usepackage[margin=2.5cm]{geometry}
\usepackage{graphicx}
\usepackage{listings}
\usepackage{xcolor}
\usepackage{hyperref}
\usepackage{booktabs}
\usepackage{longtable}
\usepackage{enumitem}
\usepackage{titlesec}
\usepackage{fancyhdr}
\usepackage{tcolorbox}
\usepackage{caption}

% ── Colour Definitions ───────────────────────────────────────────────────────
\definecolor{codebg}{HTML}{F5F5F5}
\definecolor{codeframe}{HTML}{CCCCCC}
\definecolor{javakw}{HTML}{7F0055}
\definecolor{javastr}{HTML}{2A00FF}
\definecolor{javacmt}{HTML}{3F7F5F}
\definecolor{xmltag}{HTML}{000080}
\definecolor{xmlattr}{HTML}{7F007F}
\definecolor{yamlkey}{HTML}{007F00}
\definecolor{sliitblue}{HTML}{003366}
\definecolor{accentgreen}{HTML}{28A745}
\definecolor{accentorange}{HTML}{FD7E14}

% ── Listings: Java Style ─────────────────────────────────────────────────────
\lstdefinestyle{javastyle}{
    language=Java,
    backgroundcolor=\color{codebg},
    frame=single,
    rulecolor=\color{codeframe},
    basicstyle=\ttfamily\footnotesize,
    keywordstyle=\color{javakw}\bfseries,
    stringstyle=\color{javastr},
    commentstyle=\color{javacmt}\itshape,
    numbers=left,
    numberstyle=\tiny\color{gray},
    numbersep=8pt,
    breaklines=true,
    showstringspaces=false,
    tabsize=4,
    captionpos=b,
    morekeywords={var, record, yield, sealed, permits, when},
    literate={→}{$\rightarrow$}1
}

% ── Listings: XML Style ──────────────────────────────────────────────────────
\lstdefinestyle{xmlstyle}{
    language=XML,
    backgroundcolor=\color{codebg},
    frame=single,
    rulecolor=\color{codeframe},
    basicstyle=\ttfamily\footnotesize,
    keywordstyle=\color{xmltag}\bfseries,
    stringstyle=\color{xmlattr},
    commentstyle=\color{javacmt}\itshape,
    numbers=left,
    numberstyle=\tiny\color{gray},
    numbersep=8pt,
    breaklines=true,
    showstringspaces=false,
    tabsize=4,
    captionpos=b,
    morekeywords={suite, test, classes, class}
}

% ── Listings: YAML Style ─────────────────────────────────────────────────────
\lstdefinestyle{yamlstyle}{
    basicstyle=\ttfamily\footnotesize,
    backgroundcolor=\color{codebg},
    frame=single,
    rulecolor=\color{codeframe},
    numbers=left,
    numberstyle=\tiny\color{gray},
    numbersep=8pt,
    breaklines=true,
    showstringspaces=false,
    tabsize=2,
    captionpos=b,
    keywordstyle=\color{yamlkey}\bfseries,
    commentstyle=\color{javacmt}\itshape,
    morecomment=[l]{\#}
}

% ── Listings: Gherkin Style ───────────────────────────────────────────────────
\lstdefinestyle{gherkinstyle}{
    basicstyle=\ttfamily\footnotesize,
    backgroundcolor=\color{codebg},
    frame=single,
    rulecolor=\color{codeframe},
    numbers=left,
    numberstyle=\tiny\color{gray},
    numbersep=8pt,
    breaklines=true,
    showstringspaces=false,
    tabsize=2,
    captionpos=b,
    keywordstyle=\color{javakw}\bfseries,
    stringstyle=\color{javastr},
    commentstyle=\color{javacmt}\itshape,
    morekeywords={Feature, Scenario, Given, When, Then, And, But, Background, Examples},
    morecomment=[l]{\#}
}

% ── Listings: Bash Style ─────────────────────────────────────────────────────
\lstdefinestyle{bashstyle}{
    basicstyle=\ttfamily\footnotesize,
    backgroundcolor=\color{codebg},
    frame=single,
    rulecolor=\color{codeframe},
    numbers=none,
    breaklines=true,
    showstringspaces=false,
    tabsize=2,
    captionpos=b,
    keywordstyle=\color{javakw}\bfseries,
    commentstyle=\color{javacmt}\itshape,
    morecomment=[l]{\#}
}

% ── Page Style ────────────────────────────────────────────────────────────────
\pagestyle{fancy}
\fancyhf{}
\fancyhead[L]{\footnotesize SE3010 -- Software Engineering Process \& Quality Management}
\fancyhead[R]{\footnotesize Assignment 01}
\fancyfoot[C]{\thepage}
\renewcommand{\headrulewidth}{0.4pt}

% ── tcolorbox styles ─────────────────────────────────────────────────────────
\tcbset{
    noteboxstyle/.style={
        colback=blue!5, colframe=sliitblue, fonttitle=\bfseries,
        boxrule=0.5pt, arc=2pt, left=6pt, right=6pt, top=4pt, bottom=4pt
    }
}

\hypersetup{
    colorlinks=true,
    linkcolor=sliitblue,
    urlcolor=sliitblue,
    citecolor=sliitblue
}

% ══════════════════════════════════════════════════════════════════════════════
\begin{document}

% ── Title Page ────────────────────────────────────────────────────────────────
\begin{titlepage}
    \centering
    \vspace*{2cm}
    {\Huge\bfseries\color{sliitblue} SE3010 -- Software Engineering\\Process \& Quality Management\par}
    \vspace{1cm}
    {\LARGE Assignment 01\par}
    \vspace{0.5cm}
    {\Large\itshape Part 4: Demonstration Report\par}
    \vspace{2cm}
    {\Large\bfseries Testing Tool: TestNG + Cucumber BDD\par}
    \vspace{0.3cm}
    {\large Student Management REST API\par}
    \vspace{2cm}
    \begin{tabular}{ll}
        \toprule
        \textbf{Role} & \textbf{Contribution} \\
        \midrule
        Member 1 & TestNG Assertions (Controller Integration Tests) \\
        Member 2 & TestNG Fixtures, DataProvider \& Mocking \\
        Member 3 & Cucumber BDD Step Definitions \\
        Member 4 & Cucumber Runner, Test Reporting \& CI/CD \\
        \bottomrule
    \end{tabular}
    \vspace{2cm}
    {\large Sri Lanka Institute of Information Technology\par}
    \vfill
\end{titlepage}

% ── Table of Contents ─────────────────────────────────────────────────────────
\tableofcontents
\newpage

% ══════════════════════════════════════════════════════════════════════════════
\section{Introduction}
% ══════════════════════════════════════════════════════════════════════════════

This report demonstrates the complete testing workflow for a \textbf{Student Management REST API} built with Spring Boot 3.4.3 and Java 17. The project uses \textbf{TestNG 7.10.2} as the primary test framework (JUnit is fully excluded), \textbf{Cucumber 7.18.0} for Behaviour-Driven Development (BDD), and \textbf{Mockito 5.12.0} for mocking. All tests are orchestrated through a single \texttt{testng.xml} suite, executed via Maven Surefire, and automated through a GitHub Actions CI/CD pipeline.

\subsection{Technology Stack}

\begin{table}[h!]
\centering
\begin{tabular}{lll}
    \toprule
    \textbf{Layer} & \textbf{Technology} & \textbf{Version} \\
    \midrule
    Backend Framework & Spring Boot & 3.4.3 \\
    Language & Java & 17 \\
    Test Runner & TestNG & 7.10.2 \\
    BDD Framework & Cucumber (java, testng, spring) & 7.18.0 \\
    Mocking & Mockito Core (standalone) & 5.12.0 \\
    Test Database & H2 (in-memory) & Runtime \\
    Production Database & PostgreSQL (Neon Cloud) & Runtime \\
    Build Tool & Maven + Surefire 3.2.5 & -- \\
    CI/CD & GitHub Actions & -- \\
    ORM & Spring Data JPA + Lombok & -- \\
    \bottomrule
\end{tabular}
\caption{Technology stack used in the project}
\end{table}

\subsection{Project Structure Overview}

\begin{lstlisting}[style=bashstyle, caption={Project directory structure}]
se3010-testing/
|-- pom.xml                       # Maven config
|-- testng.xml                    # TestNG suite definition
|-- .github/workflows/test.yml    # CI/CD pipeline
|-- src/main/java/com/sliit/se3010_testing/
|   |-- Se3010TestingApplication.java
|   |-- model/Student.java
|   |-- repository/StudentRepository.java
|   |-- service/StudentService.java
|   |-- controller/StudentController.java
|-- src/main/resources/
|   |-- application.properties    # PostgreSQL config
|-- src/test/java/
|   |-- tests/StudentControllerTest.java   # Member 1
|   |-- tests/StudentServiceTest.java      # Member 2
|   |-- tests/CucumberRunnerTest.java      # Member 4
|   |-- steps/StudentSteps.java            # Member 3
|-- src/test/resources/
    |-- application-test.properties        # H2 config
    |-- cucumber.properties
    |-- features/student_registration.feature  # BDD scenarios
\end{lstlisting}

% ══════════════════════════════════════════════════════════════════════════════
\section{Application Under Test}
% ══════════════════════════════════════════════════════════════════════════════

\subsection{Domain Model -- \texttt{Student.java}}

The \texttt{Student} entity uses Lombok annotations to eliminate boilerplate code. The \texttt{@Builder} pattern enables fluent object creation in tests.

\begin{lstlisting}[style=javastyle, caption={Student entity model}]
@Entity
@Data
@NoArgsConstructor
@AllArgsConstructor
@Builder
public class Student {
    @Id
    @GeneratedValue(strategy = GenerationType.IDENTITY)
    private Long id;
    
    private String name;
    private String email;
    private int age;
    private String course;
}
\end{lstlisting}

\subsection{Repository -- \texttt{StudentRepository.java}}

Extends \texttt{JpaRepository} with a custom finder method used for duplicate email detection.

\begin{lstlisting}[style=javastyle, caption={Student repository interface}]
@Repository
public interface StudentRepository extends JpaRepository<Student, Long> {
    Optional<Student> findByEmail(String email);
}
\end{lstlisting}

\subsection{Service Layer -- \texttt{StudentService.java}}

Contains business logic including a custom \texttt{DuplicateEmailException}. This is the class tested by Member~2 using Mockito.

\begin{lstlisting}[style=javastyle, caption={Student service with business logic}]
@Service
@RequiredArgsConstructor
public class StudentService {

    private final StudentRepository studentRepository;

    public static class DuplicateEmailException extends RuntimeException {
        public DuplicateEmailException(String email) {
            super("A student with email '" + email + "' already exists.");
        }
    }

    public Student createStudent(Student student) {
        studentRepository.findByEmail(student.getEmail())
            .ifPresent(existing -> {
                throw new DuplicateEmailException(student.getEmail());
            });
        return studentRepository.save(student);
    }

    public List<Student> getAllStudents() {
        return studentRepository.findAll();
    }

    public Optional<Student> getStudentById(Long id) {
        return studentRepository.findById(id);
    }

    public Student updateStudent(Long id, Student studentDetails) {
        return studentRepository.findById(id).map(student -> {
            student.setName(studentDetails.getName());
            student.setEmail(studentDetails.getEmail());
            student.setAge(studentDetails.getAge());
            student.setCourse(studentDetails.getCourse());
            return studentRepository.save(student);
        }).orElseThrow(() ->
            new RuntimeException("Student not found with id " + id));
    }

    public void deleteStudent(Long id) {
        studentRepository.deleteById(id);
    }
}
\end{lstlisting}

\subsection{REST Controller -- \texttt{StudentController.java}}

Exposes CRUD endpoints at \texttt{/api/students}. This is the class tested by Member~1 using MockMvc.

\begin{lstlisting}[style=javastyle, caption={REST controller endpoints}]
@RestController
@RequestMapping("/api/students")
@RequiredArgsConstructor
@CrossOrigin(origins = "*")
public class StudentController {

    private final StudentService studentService;

    @PostMapping
    public ResponseEntity<?> createStudent(@RequestBody Student student) {
        try {
            Student created = studentService.createStudent(student);
            return new ResponseEntity<>(created, HttpStatus.CREATED);
        } catch (StudentService.DuplicateEmailException e) {
            return ResponseEntity.status(HttpStatus.CONFLICT)
                                 .body(e.getMessage());
        }
    }

    @GetMapping
    public ResponseEntity<List<Student>> getAllStudents() {
        return ResponseEntity.ok(studentService.getAllStudents());
    }

    @GetMapping("/{id}")
    public ResponseEntity<Student> getStudentById(@PathVariable Long id) {
        return studentService.getStudentById(id)
                .map(ResponseEntity::ok)
                .orElse(ResponseEntity.notFound().build());
    }

    @PutMapping("/{id}")
    public ResponseEntity<Student> updateStudent(
            @PathVariable Long id, @RequestBody Student details) {
        try {
            return ResponseEntity.ok(
                studentService.updateStudent(id, details));
        } catch (RuntimeException e) {
            return ResponseEntity.notFound().build();
        }
    }

    @DeleteMapping("/{id}")
    public ResponseEntity<Void> deleteStudent(@PathVariable Long id) {
        studentService.deleteStudent(id);
        return ResponseEntity.noContent().build();
    }
}
\end{lstlisting}

\subsection{REST API Endpoints}

\begin{table}[h!]
\centering
\begin{tabular}{llll}
    \toprule
    \textbf{Method} & \textbf{Endpoint} & \textbf{Success} & \textbf{Error} \\
    \midrule
    POST & \texttt{/api/students} & 201 Created & 409 Conflict \\
    GET & \texttt{/api/students} & 200 OK & -- \\
    GET & \texttt{/api/students/\{id\}} & 200 OK & 404 Not Found \\
    PUT & \texttt{/api/students/\{id\}} & 200 OK & 404 Not Found \\
    DELETE & \texttt{/api/students/\{id\}} & 204 No Content & -- \\
    \bottomrule
\end{tabular}
\caption{REST API endpoint summary}
\end{table}

% ══════════════════════════════════════════════════════════════════════════════
\section{TestNG Suite Configuration}
% ══════════════════════════════════════════════════════════════════════════════

All tests are organised in a single TestNG XML suite that maps each member's contribution to a named test block.

\begin{lstlisting}[style=xmlstyle, caption={\texttt{testng.xml} -- Test suite definition}]
<?xml version="1.0" encoding="UTF-8"?>
<!DOCTYPE suite SYSTEM "https://testng.org/testng-1.0.dtd">

<suite name="Student API Test Suite" verbose="2">

    <!-- Member 1: TestNG Assertions (MockMvc controller tests) -->
    <test name="Member1 - Controller Assertions">
        <classes>
            <class name="tests.StudentControllerTest"/>
        </classes>
    </test>

    <!-- Member 2: TestNG Fixtures + DataProvider (service unit tests) -->
    <test name="Member2 - Service Fixtures">
        <classes>
            <class name="tests.StudentServiceTest"/>
        </classes>
    </test>

    <!-- Member 3 + 4: Cucumber BDD via TestNG runner -->
    <test name="Member3+4 - Cucumber BDD">
        <classes>
            <class name="tests.CucumberRunnerTest"/>
        </classes>
    </test>

</suite>
\end{lstlisting}

\begin{tcolorbox}[noteboxstyle, title=Key Configuration Notes]
\begin{itemize}[nosep]
    \item \textbf{verbose="2"} -- Provides detailed output including test method names and status.
    \item Each \texttt{<test>} block runs independently with its own Spring context if needed.
    \item Maven Surefire is configured in \texttt{pom.xml} to use this file via\\
    \texttt{<suiteXmlFile>testng.xml</suiteXmlFile>}.
    \item JUnit is \textbf{completely excluded} from \texttt{pom.xml} using six explicit exclusions.
\end{itemize}
\end{tcolorbox}

% ══════════════════════════════════════════════════════════════════════════════
\section{Test Database Configuration}
% ══════════════════════════════════════════════════════════════════════════════

Tests use \texttt{@ActiveProfiles("test")} which loads a separate properties file with an H2 in-memory database, ensuring complete isolation from the production PostgreSQL database.

\begin{lstlisting}[style=bashstyle, caption={\texttt{application-test.properties} -- H2 test database}]
spring.datasource.url=jdbc:h2:mem:testdb;DB_CLOSE_DELAY=-1
spring.datasource.driver-class-name=org.h2.Driver
spring.datasource.username=sa
spring.datasource.password=
spring.jpa.database-platform=org.hibernate.dialect.H2Dialect
spring.jpa.hibernate.ddl-auto=create-drop
spring.jpa.show-sql=true
\end{lstlisting}

\begin{tcolorbox}[noteboxstyle, title=Why H2 for Tests?]
\begin{itemize}[nosep]
    \item \textbf{create-drop} -- Schema is created fresh before tests and dropped after, ensuring no leftover data.
    \item \textbf{In-memory} -- No disk I/O means faster test execution.
    \item \textbf{Isolation} -- Tests never touch the production Neon PostgreSQL database.
\end{itemize}
\end{tcolorbox}

% ══════════════════════════════════════════════════════════════════════════════
\section{Member 1: TestNG Assertions}
\label{sec:member1}
% ══════════════════════════════════════════════════════════════════════════════

\subsection{Feature Demonstrated}
\textbf{TestNG Assertions} -- Using \texttt{Assert.assertEquals()}, \texttt{Assert.assertFalse()}, and \texttt{Assert.assertNotNull()} to validate HTTP responses from a Spring Boot REST controller via MockMvc.

\subsection{Key Concepts}
\begin{itemize}
    \item \textbf{AbstractTestNGSpringContextTests} -- Bridge class that allows TestNG (instead of JUnit) to manage the Spring application context.
    \item \textbf{@SpringBootTest + @AutoConfigureMockMvc} -- Starts the full Spring context with a mock servlet environment.
    \item \textbf{@BeforeMethod} -- TestNG fixture annotation that clears the database before each test method.
    \item \textbf{MockMvc} -- Spring's testing utility to perform HTTP requests without starting a real server.
\end{itemize}

\subsection{Complete Test Class}

\begin{lstlisting}[style=javastyle, caption={\texttt{StudentControllerTest.java} -- Member 1's assertion tests}]
@SpringBootTest(classes = Se3010TestingApplication.class,
        webEnvironment = SpringBootTest.WebEnvironment.MOCK)
@AutoConfigureMockMvc
@ActiveProfiles("test")
public class StudentControllerTest
        extends AbstractTestNGSpringContextTests {

    @Autowired private MockMvc mockMvc;
    @Autowired private StudentRepository studentRepository;
    @Autowired private ObjectMapper objectMapper;

    @BeforeMethod
    public void cleanUp() {
        studentRepository.deleteAll();
    }

    // Test 1: POST /api/students -> 201 Created
    @Test
    public void testCreateStudent_Returns201AndCorrectBody()
            throws Exception {
        Student student = Student.builder()
            .name("Alice")
            .email("alice@test.com")
            .age(22)
            .course("SE3010")
            .build();

        String body = mockMvc.perform(post("/api/students")
                .contentType(MediaType.APPLICATION_JSON)
                .content(objectMapper.writeValueAsString(student)))
            .andExpect(status().isCreated())
            .andExpect(jsonPath("$.name").value("Alice"))
            .andExpect(jsonPath("$.email").value("alice@test.com"))
            .andReturn().getResponse().getContentAsString();

        Assert.assertFalse(body.isEmpty(),
            "Response body must not be empty");
    }

    // Test 2: GET /api/students/{id} -> 200 OK
    @Test
    public void testGetStudent_Returns200() throws Exception {
        Student saved = studentRepository.save(
            Student.builder()
                .name("Bob").email("bob@test.com")
                .age(23).course("SE3020")
                .build());

        String body = mockMvc.perform(
                get("/api/students/{id}", saved.getId()))
            .andExpect(status().isOk())
            .andExpect(jsonPath("$.name").value("Bob"))
            .andReturn().getResponse().getContentAsString();

        Assert.assertNotNull(body,
            "Response body should not be null");
    }

    // Test 3: GET /api/students/99999 -> 404 Not Found
    @Test
    public void testGetStudent_NotFound_Returns404()
            throws Exception {
        int statusCode = mockMvc.perform(
                get("/api/students/{id}", 99999L))
            .andReturn().getResponse().getStatus();

        Assert.assertEquals(statusCode, 404,
            "Expected HTTP 404 for missing student");
    }

    // Test 4: POST with malformed JSON -> 400 Bad Request
    @Test
    public void testCreateStudent_MalformedJson_Returns400()
            throws Exception {
        int statusCode = mockMvc.perform(post("/api/students")
                .contentType(MediaType.APPLICATION_JSON)
                .content("{ invalid json }"))
            .andReturn().getResponse().getStatus();

        Assert.assertEquals(statusCode, 400,
            "Expected HTTP 400 for malformed JSON");
    }
}
\end{lstlisting}

\subsection{Test Cases Summary}

\begin{table}[h!]
\centering
\begin{tabular}{clllc}
    \toprule
    \textbf{\#} & \textbf{Test Method} & \textbf{HTTP} & \textbf{Assertion Used} & \textbf{Expected} \\
    \midrule
    1 & \texttt{testCreateStudent\_Returns201...} & POST & \texttt{Assert.assertFalse} & 201 \\
    2 & \texttt{testGetStudent\_Returns200} & GET & \texttt{Assert.assertNotNull} & 200 \\
    3 & \texttt{testGetStudent\_NotFound...} & GET & \texttt{Assert.assertEquals} & 404 \\
    4 & \texttt{testCreateStudent\_Malformed...} & POST & \texttt{Assert.assertEquals} & 400 \\
    \bottomrule
\end{tabular}
\caption{Member 1 -- Test cases for controller assertions}
\end{table}

\subsection{Code Explanation}

\begin{enumerate}
    \item \textbf{Line 1--3}: \texttt{@SpringBootTest} loads the full application context in \texttt{MOCK} mode (no real HTTP server). \texttt{@AutoConfigureMockMvc} auto-configures \texttt{MockMvc}.
    \item \textbf{Line 5--6}: The class extends \texttt{AbstractTestNGSpringContextTests} -- this is the \textbf{key difference from JUnit}. Without this, Spring context would not initialise under TestNG.
    \item \textbf{Line 14--16}: \texttt{@BeforeMethod} is TestNG's equivalent of JUnit's \texttt{@BeforeEach}. It runs before \textit{every} test method.
    \item \textbf{Lines 21--35}: Uses MockMvc's fluent API: \texttt{perform()} $\rightarrow$ \texttt{andExpect()} $\rightarrow$ \texttt{andReturn()} chain. JSON path assertions validate response body structure.
    \item \textbf{Line 35}: \texttt{Assert.assertFalse(body.isEmpty())} -- TestNG assertion confirming the response is not empty.
\end{enumerate}

% ══════════════════════════════════════════════════════════════════════════════
\section{Member 2: TestNG Fixtures, DataProvider \& Mocking}
\label{sec:member2}
% ══════════════════════════════════════════════════════════════════════════════

\subsection{Features Demonstrated}
\begin{enumerate}
    \item \textbf{TestNG Fixtures} -- \texttt{@BeforeMethod} and \texttt{@AfterMethod} for setup/teardown.
    \item \textbf{TestNG DataProvider} -- \texttt{@DataProvider} for parameterised test execution.
    \item \textbf{Mockito Mocking} -- \texttt{mock()}, \texttt{when().thenReturn()}, \texttt{verify()} without JUnit integration.
\end{enumerate}

\subsection{Key Concepts}
\begin{itemize}
    \item \textbf{Pure unit test} -- No Spring context loaded. The repository is fully mocked.
    \item \textbf{@BeforeMethod} runs before every test \textit{and} every DataProvider row.
    \item \textbf{@AfterMethod} calls \texttt{verifyNoMoreInteractions()} to catch unexpected mock calls.
    \item \textbf{@DataProvider} returns an \texttt{Object[][]} matrix. Each row becomes a separate test invocation.
\end{itemize}

\subsection{Complete Test Class}

\begin{lstlisting}[style=javastyle, caption={\texttt{StudentServiceTest.java} -- Member 2's fixture and mocking tests}]
public class StudentServiceTest {

    private StudentRepository studentRepository;
    private StudentService studentService;

    // FIXTURE: Fresh mocks before every test / DataProvider row
    @BeforeMethod
    public void setUp() {
        studentRepository = mock(StudentRepository.class);
        studentService = new StudentService(studentRepository);
        System.out.println(
            "[FIXTURE] Fresh mocks created - test starting");
    }

    // FIXTURE: Verify no unexpected interactions after each test
    @AfterMethod
    public void tearDown() {
        verifyNoMoreInteractions(studentRepository);
        System.out.println(
            "[FIXTURE] Verified - test finished");
    }

    // DATAPROVIDER: Multiple student records
    @DataProvider(name = "studentData")
    public Object[][] studentDataProvider() {
        return new Object[][] {
            {"Alice", "alice@test.com", 22, "SE3010"},
            {"Bob",   "bob@test.com",   25, "SE3020"},
            {"Carol", "carol@test.com", 20, "SE3030"},
        };
    }

    // Test 1: createStudent with parameterised data (runs 3x)
    @Test(dataProvider = "studentData")
    public void testCreateStudent(String name, String email,
            int age, String course) {
        Student input = Student.builder()
            .name(name).email(email)
            .age(age).course(course).build();
        Student saved = Student.builder()
            .id(1L).name(name).email(email)
            .age(age).course(course).build();

        when(studentRepository.findByEmail(email))
            .thenReturn(Optional.empty());
        when(studentRepository.save(any(Student.class)))
            .thenReturn(saved);

        Student result = studentService.createStudent(input);

        Assert.assertNotNull(result.getId(),
            "Saved student must have an ID");
        Assert.assertEquals(result.getName(), name,
            "Name should match");
        Assert.assertEquals(result.getEmail(), email,
            "Email should match");
        verify(studentRepository).findByEmail(email);
        verify(studentRepository).save(any(Student.class));
    }

    // Test 2: getStudentById - found
    @Test
    public void testGetStudentById_Found() {
        Student student = Student.builder()
            .id(1L).name("Dave").email("dave@test.com")
            .age(21).course("SE3040").build();
        when(studentRepository.findById(1L))
            .thenReturn(Optional.of(student));

        Optional<Student> result =
            studentService.getStudentById(1L);

        Assert.assertTrue(result.isPresent(),
            "Student should be present");
        Assert.assertEquals(result.get().getName(), "Dave");
        verify(studentRepository).findById(1L);
    }

    // Test 3: getStudentById - not found
    @Test
    public void testGetStudentById_NotFound() {
        when(studentRepository.findById(99L))
            .thenReturn(Optional.empty());

        Optional<Student> result =
            studentService.getStudentById(99L);

        Assert.assertFalse(result.isPresent(),
            "Student should not be present");
        verify(studentRepository).findById(99L);
    }

    // Test 4: getAllStudents returns list
    @Test
    public void testGetAllStudents_ReturnsList() {
        List<Student> list = List.of(
            Student.builder().id(1L).name("Eve")
                .email("eve@test.com").age(23)
                .course("SE3050").build());
        when(studentRepository.findAll()).thenReturn(list);

        List<Student> result =
            studentService.getAllStudents();

        Assert.assertEquals(result.size(), 1,
            "Should return exactly one student");
        Assert.assertEquals(result.get(0).getName(), "Eve");
        verify(studentRepository).findAll();
    }
}
\end{lstlisting}

\subsection{Test Cases Summary}

\begin{table}[h!]
\centering
\begin{tabular}{cllll}
    \toprule
    \textbf{\#} & \textbf{Test Method} & \textbf{Data} & \textbf{Mock Setup} & \textbf{Runs} \\
    \midrule
    1 & \texttt{testCreateStudent} & DataProvider & \texttt{findByEmail} + \texttt{save} & 3$\times$ \\
    2 & \texttt{testGetStudentById\_Found} & Inline & \texttt{findById} $\rightarrow$ present & 1$\times$ \\
    3 & \texttt{testGetStudentById\_NotFound} & Inline & \texttt{findById} $\rightarrow$ empty & 1$\times$ \\
    4 & \texttt{testGetAllStudents\_ReturnsList} & Inline & \texttt{findAll} $\rightarrow$ list & 1$\times$ \\
    \bottomrule
\end{tabular}
\caption{Member 2 -- Test cases (6 total executions: 3 parameterised + 3 individual)}
\end{table}

\subsection{Code Explanation}

\begin{enumerate}
    \item \textbf{@BeforeMethod (line 7--13)}: Creates a fresh \texttt{mock(StudentRepository.class)} and injects it into a new \texttt{StudentService} instance. This runs 6 times total (once per test execution).
    \item \textbf{@AfterMethod (line 16--21)}: \texttt{verifyNoMoreInteractions()} ensures no unexpected repository calls happened. If a test calls \texttt{findById} but forgets to \texttt{verify()} it, this will fail.
    \item \textbf{@DataProvider (line 24--31)}: Returns 3 rows of \texttt{\{name, email, age, course\}}. TestNG maps each row to the parameters of \texttt{testCreateStudent()}.
    \item \textbf{when().thenReturn() (line 44--47)}: Configures mock behaviour -- tells the mock what to return when specific methods are called.
    \item \textbf{verify() (line 56--57)}: Confirms that the service actually called the expected repository methods. Combined with \texttt{verifyNoMoreInteractions()}, this provides full interaction verification.
\end{enumerate}

% ══════════════════════════════════════════════════════════════════════════════
\section{Member 3: Cucumber BDD Step Definitions}
\label{sec:member3}
% ══════════════════════════════════════════════════════════════════════════════

\subsection{Feature Demonstrated}
\textbf{Behaviour-Driven Development (BDD)} -- Writing test scenarios in plain English (Gherkin syntax) and implementing step definitions in Java.

\subsection{Key Concepts}
\begin{itemize}
    \item \textbf{Gherkin} -- A business-readable language using \texttt{Feature}, \texttt{Scenario}, \texttt{Given}, \texttt{When}, \texttt{Then} keywords.
    \item \textbf{@CucumberContextConfiguration} -- Marks the step definition class as the Spring context configuration source for Cucumber.
    \item \textbf{cucumber-spring} -- Integrates Cucumber with the Spring Boot test context.
    \item \textbf{Step Definitions} -- Java methods annotated with \texttt{@Given}, \texttt{@When}, \texttt{@Then} that map to Gherkin steps.
\end{itemize}

\subsection{Feature File -- Gherkin Scenarios}

\begin{lstlisting}[style=gherkinstyle, caption={\texttt{student\_registration.feature} -- BDD scenarios in Gherkin}]
Feature: Student Registration API
  As an admin
  I want to register students via the API
  So that I can manage the student roster

  # Scenario 1: Successful registration (201)
  Scenario: Successful student registration
    Given no students exist in the system
    When I register a student with name "Alice"
         email "alice@uni.com" age 22 course "SE3010"
    Then the registration response status should be 201
    And the registered student name should be "Alice"

  # Scenario 2: Duplicate email (409)
  Scenario: Duplicate email returns conflict error
    Given a student with email "bob@uni.com" already exists
    When I register a student with name "Bob2"
         email "bob@uni.com" age 24 course "SE3020"
    Then the registration response status should be 409

  # Scenario 3: Missing required field (400)
  Scenario: Missing required field returns bad request error
    Given no students exist in the system
    When I send a registration request with malformed JSON body
    Then the registration response status should be 400
\end{lstlisting}

\subsection{Step Definitions -- Java Implementation}

\begin{lstlisting}[style=javastyle, caption={\texttt{StudentSteps.java} -- Cucumber step definitions}]
@CucumberContextConfiguration
@SpringBootTest(classes = Se3010TestingApplication.class,
        webEnvironment = SpringBootTest.WebEnvironment.MOCK)
@AutoConfigureMockMvc
@ActiveProfiles("test")
public class StudentSteps {

    @Autowired private MockMvc mockMvc;
    @Autowired private StudentRepository studentRepository;
    @Autowired private ObjectMapper objectMapper;

    private MvcResult lastResult;

    // Cucumber hook: reset DB before each scenario
    @Before
    public void resetDatabase() {
        studentRepository.deleteAll();
    }

    // GIVEN: No students exist
    @Given("no students exist in the system")
    public void no_students_exist() {
        studentRepository.deleteAll();
    }

    // GIVEN: A student with email already exists
    @Given("a student with email {string} already exists")
    public void a_student_with_email_already_exists(
            String email) {
        studentRepository.deleteAll();
        studentRepository.save(
            Student.builder()
                .name("Existing").email(email)
                .age(20).course("SE0000")
                .build());
    }

    // WHEN: Register a student
    @When("I register a student with name {string} "
        + "email {string} age {int} course {string}")
    public void i_register_a_student(String name,
            String email, int age, String course)
            throws Exception {
        Student s = Student.builder()
            .name(name).email(email)
            .age(age).course(course).build();
        lastResult = mockMvc.perform(post("/api/students")
                .contentType(MediaType.APPLICATION_JSON)
                .content(objectMapper.writeValueAsString(s)))
            .andReturn();
    }

    // WHEN: Send malformed JSON
    @When("I send a registration request with "
        + "malformed JSON body")
    public void i_send_malformed_json() throws Exception {
        lastResult = mockMvc.perform(post("/api/students")
                .contentType(MediaType.APPLICATION_JSON)
                .content("{ bad json }"))
            .andReturn();
    }

    // THEN: Check HTTP status code
    @Then("the registration response status "
        + "should be {int}")
    public void the_registration_status_should_be(
            int expectedStatus) {
        int actual = lastResult.getResponse().getStatus();
        Assert.assertEquals(actual, expectedStatus,
            "Expected HTTP " + expectedStatus
            + " but got " + actual);
    }

    // THEN: Check student name in response
    @Then("the registered student name "
        + "should be {string}")
    public void the_registered_student_name_should_be(
            String expectedName) throws Exception {
        String body = lastResult.getResponse()
            .getContentAsString();
        Student returned = objectMapper
            .readValue(body, Student.class);
        Assert.assertEquals(returned.getName(),
            expectedName, "Student name mismatch");
    }
}
\end{lstlisting}

\subsection{Test Cases Summary}

\begin{table}[h!]
\centering
\begin{tabular}{clll}
    \toprule
    \textbf{\#} & \textbf{Scenario} & \textbf{Steps} & \textbf{Expected Status} \\
    \midrule
    1 & Successful registration & Given $\rightarrow$ When $\rightarrow$ Then $\rightarrow$ And & 201 Created \\
    2 & Duplicate email conflict & Given $\rightarrow$ When $\rightarrow$ Then & 409 Conflict \\
    3 & Malformed JSON & Given $\rightarrow$ When $\rightarrow$ Then & 400 Bad Request \\
    \bottomrule
\end{tabular}
\caption{Member 3 -- Cucumber BDD scenarios}
\end{table}

\subsection{Gherkin to Java Mapping}

\begin{table}[h!]
\centering
\small
\begin{tabular}{lll}
    \toprule
    \textbf{Gherkin Step} & \textbf{Annotation} & \textbf{Java Method} \\
    \midrule
    \texttt{no students exist in the system} & \texttt{@Given} & \texttt{no\_students\_exist()} \\
    \texttt{a student with email "..." already exists} & \texttt{@Given} & \texttt{a\_student\_with\_email...} \\
    \texttt{I register a student with name "..."} & \texttt{@When} & \texttt{i\_register\_a\_student(...)} \\
    \texttt{I send ... malformed JSON body} & \texttt{@When} & \texttt{i\_send\_malformed\_json()} \\
    \texttt{the ... status should be \{int\}} & \texttt{@Then} & \texttt{the\_registration\_status...} \\
    \texttt{the ... name should be "..."} & \texttt{@Then} & \texttt{the\_registered\_student...} \\
    \bottomrule
\end{tabular}
\caption{Gherkin step to Java step definition mapping}
\end{table}

% ══════════════════════════════════════════════════════════════════════════════
\section{Member 4: Cucumber TestNG Runner, Reporting \& CI/CD}
\label{sec:member4}
% ══════════════════════════════════════════════════════════════════════════════

\subsection{Features Demonstrated}
\begin{enumerate}
    \item \textbf{Test Runner} -- Cucumber TestNG runner extending \texttt{AbstractTestNGCucumberTests}.
    \item \textbf{Test Reporting} -- Cucumber HTML/JSON reports + Maven Surefire reports.
    \item \textbf{CI/CD Pipeline} -- GitHub Actions workflow for automated testing.
\end{enumerate}

\subsection{Cucumber TestNG Runner}

\begin{lstlisting}[style=javastyle, caption={\texttt{CucumberRunnerTest.java} -- TestNG-based Cucumber runner}]
@CucumberOptions(
    features = "src/test/resources/features/"
             + "student_registration.feature",
    glue     = {"steps"},
    plugin   = {
        "pretty",
        "html:target/cucumber-reports/cucumber.html",
        "json:target/cucumber-reports/cucumber.json"
    },
    monochrome = true
)
public class CucumberRunnerTest
        extends AbstractTestNGCucumberTests {

    /**
     * Override to allow parallel scenario execution
     * (set parallel = true to enable).
     */
    @Override
    @DataProvider(parallel = false)
    public Object[][] scenarios() {
        return super.scenarios();
    }
}
\end{lstlisting}

\subsection{Code Explanation -- Runner}

\begin{enumerate}
    \item \textbf{@CucumberOptions}: Configures where to find feature files (\texttt{features}), step definitions (\texttt{glue}), and output format (\texttt{plugin}).
    \item \textbf{AbstractTestNGCucumberTests}: This is the \textbf{TestNG equivalent} of JUnit's \texttt{@RunWith(Cucumber.class)}. It bridges Cucumber into the TestNG execution lifecycle.
    \item \textbf{plugin -- ``pretty''}: Prints Gherkin steps with colour in the console during execution.
    \item \textbf{plugin -- ``html:...''}: Generates a standalone HTML report at \texttt{target/cucumber-reports/cucumber.html}.
    \item \textbf{plugin -- ``json:...''}: Generates a JSON report for integration with external tools.
    \item \textbf{monochrome = true}: Removes ANSI colour codes for cleaner log output.
    \item \textbf{@DataProvider(parallel = false)}: Each Cucumber scenario becomes a TestNG DataProvider row. Setting \texttt{parallel = true} would run scenarios concurrently.
\end{enumerate}

\subsection{Test Reports Generated}

\begin{table}[h!]
\centering
\begin{tabular}{lll}
    \toprule
    \textbf{Report} & \textbf{Location} & \textbf{Format} \\
    \midrule
    Cucumber HTML & \texttt{target/cucumber-reports/cucumber.html} & Standalone HTML \\
    Cucumber JSON & \texttt{target/cucumber-reports/cucumber.json} & JSON \\
    TestNG Default & \texttt{target/surefire-reports/} & XML + HTML \\
    Surefire Report & \texttt{target/site/surefire-report.html} & HTML \\
    \bottomrule
\end{tabular}
\caption{Generated test reports and their locations}
\end{table}

\subsection{Maven Surefire Report Plugin}

The following plugin in \texttt{pom.xml} generates an HTML report automatically during the \texttt{test} phase:

\begin{lstlisting}[style=xmlstyle, caption={Surefire report plugin configuration in \texttt{pom.xml}}]
<!-- Surefire Report Plugin to generate HTML report -->
<plugin>
    <groupId>org.apache.maven.plugins</groupId>
    <artifactId>maven-surefire-report-plugin</artifactId>
    <version>3.2.5</version>
    <executions>
        <execution>
            <phase>test</phase>
            <goals>
                <goal>report</goal>
            </goals>
        </execution>
    </executions>
</plugin>
\end{lstlisting}

\subsection{CI/CD Pipeline -- GitHub Actions}

\begin{lstlisting}[style=yamlstyle, caption={\texttt{.github/workflows/test.yml} -- CI/CD pipeline}]
name: Spring Boot CI/CD Pipeline

on:
  push:
    branches: [ "main", "master" ]
  pull_request:
    branches: [ "main", "master" ]

jobs:
  build:
    runs-on: ubuntu-latest

    steps:
    - uses: actions/checkout@v3

    - name: Set up JDK 17
      uses: actions/setup-java@v3
      with:
        java-version: '17'
        distribution: 'temurin'
        cache: 'maven'

    - name: Build and run tests with Maven
      run: mvn clean test
      working-directory: .

    - name: Publish Test Report
      if: success() || failure()
      uses: actions/upload-artifact@v3
      with:
        name: surefire-test-reports
        path: target/surefire-reports/
\end{lstlisting}

\subsection{CI/CD Pipeline Explanation}

\begin{enumerate}
    \item \textbf{Trigger (lines 3--7)}: The pipeline runs automatically on every \texttt{push} or \texttt{pull\_request} to the \texttt{main} or \texttt{master} branch.
    
    \item \textbf{Checkout (line 14)}: \texttt{actions/checkout@v3} clones the repository into the runner.
    
    \item \textbf{JDK Setup (lines 16--21)}: \texttt{actions/setup-java@v3} installs JDK~17 (Temurin distribution) and caches Maven dependencies for faster subsequent runs.
    
    \item \textbf{Test Execution (lines 23--25)}: \texttt{mvn clean test} compiles the project and runs the entire \texttt{testng.xml} suite, which includes:
    \begin{itemize}[nosep]
        \item Member 1's controller assertion tests
        \item Member 2's service mock/fixture tests
        \item Member 3+4's Cucumber BDD scenarios
    \end{itemize}
    
    \item \textbf{Report Upload (lines 27--32)}: \texttt{actions/upload-artifact@v3} uploads the \texttt{target/surefire-reports/} directory as a downloadable artifact in GitHub. The \texttt{if: success() || failure()} condition ensures reports are uploaded \textbf{even when tests fail}, making it easy to diagnose failures.
\end{enumerate}

\begin{tcolorbox}[noteboxstyle, title=CI/CD Workflow Summary]
\centering
\texttt{git push} $\rightarrow$ GitHub Actions $\rightarrow$ JDK 17 Setup $\rightarrow$ \texttt{mvn clean test} $\rightarrow$ Upload Reports
\end{tcolorbox}

% ══════════════════════════════════════════════════════════════════════════════
\section{How to Run the Tests}
% ══════════════════════════════════════════════════════════════════════════════

\subsection{Prerequisites}
\begin{itemize}
    \item Java 17 (JDK) installed
    \item Apache Maven 3.8+ installed
    \item Git (for CI/CD)
\end{itemize}

\subsection{Running All Tests Locally}

\begin{lstlisting}[style=bashstyle, caption={Run the complete test suite}]
# Navigate to the backend directory
cd se3010-testing

# Run all tests via TestNG suite
mvn clean test

# This executes testng.xml which runs:
#   - Member 1: StudentControllerTest (4 tests)
#   - Member 2: StudentServiceTest    (6 executions)
#   - Member 3+4: CucumberRunnerTest  (3 scenarios)
\end{lstlisting}

\subsection{Viewing Test Reports}

\begin{lstlisting}[style=bashstyle, caption={Locate and open test reports after execution}]
# TestNG / Surefire XML reports
ls target/surefire-reports/

# Cucumber HTML report (open in browser)
open target/cucumber-reports/cucumber.html

# Maven Surefire HTML report (if generated)
open target/site/surefire-report.html
\end{lstlisting}

\subsection{Running the Backend Server}

\begin{lstlisting}[style=bashstyle, caption={Start the Spring Boot application}]
cd se3010-testing
mvn spring-boot:run
# Server starts at http://localhost:8080
# API available at http://localhost:8080/api/students
\end{lstlisting}

\subsection{Running the Frontend}

\begin{lstlisting}[style=bashstyle, caption={Start the React frontend}]
cd frontend
npm install
npm run dev
# Frontend starts at http://localhost:5173
\end{lstlisting}

% ══════════════════════════════════════════════════════════════════════════════
\section{Summary of Demonstrated Features}
% ══════════════════════════════════════════════════════════════════════════════

\begin{table}[h!]
\centering
\begin{tabular}{clll}
    \toprule
    \textbf{Member} & \textbf{Feature} & \textbf{Key File} & \textbf{Tests} \\
    \midrule
    1 & Assertions & \texttt{StudentControllerTest.java} & 4 \\
    2 & Fixtures + Mocking & \texttt{StudentServiceTest.java} & 6 \\
    3 & BDD Syntax & \texttt{StudentSteps.java} + \texttt{.feature} & 3 \\
    4 & Reporting + CI/CD & \texttt{CucumberRunnerTest.java} + \texttt{test.yml} & -- \\
    \midrule
    \multicolumn{3}{r}{\textbf{Total Test Executions}} & \textbf{13} \\
    \bottomrule
\end{tabular}
\caption{Summary of all demonstrated features across 4 group members}
\end{table}

\begin{tcolorbox}[noteboxstyle, title=Key Takeaway]
The project demonstrates \textbf{6 distinct testing features} (Assertions, Fixtures, Mocking, DataProvider, BDD, Reporting, and CI/CD) using \textbf{TestNG as the sole test runner} -- JUnit is completely excluded. All tests are executable via a single \texttt{mvn clean test} command and automated through GitHub Actions.
\end{tcolorbox}

\end{document}
